\section{Incomplete and regularized beta functions}
\label{sec:incomplete}

\subsection{Effective definition}

Let $x$ be a real name with $0\le x\le1$.  For $\Rea a>0$, define
\[
 \B_x(a,b)=\int_0^x t^{a-1}(1-t)^{b-1}\,dt.
\]
If $x\le1-\delta$ with $\delta>0$, only the lower endpoint is improper and the compact derivative bound uses the margin $1-t\ge\delta$.  At $x=1$ use the complete-beta construction.  If the sign or endpoint status of a general name $x$ is not decidable, the theorem is stated under one of these positive certificates; no exact-real case split is made.

The regularized incomplete beta function is
\[
 \I_x(a,b)=\frac{\B_x(a,b)}{\B(a,b)}
\]
for positive real parameters, where \cref{thm:basic-laws} supplies a positive lower bound on the denominator over a certified compact domain.

\subsection{Basic identities}

\begin{theorem}[Endpoints, complement, and elementary cases]\label{thm:incomplete-basic}
For positive real parameters and $0\le x\le1$,
\begin{align*}
 \B_0(a,b)&=0,& \B_1(a,b)&=\B(a,b),\\
 \B_x(a,b)+\B_{1-x}(b,a)&=\B(a,b),\\
 \B_x(a,1)&=\frac{x^a}{a},&
 \B_x(1,b)&=\frac{1-(1-x)^b}{b}.
\end{align*}
Consequently
\begin{align*}
 \I_0(a,b)&=0,& \I_1(a,b)&=1,\\
 \I_x(a,b)&=1-\I_{1-x}(b,a),\\
 \I_x(a,1)&=x^a,&
 \I_x(1,b)&=1-(1-x)^b.
\end{align*}
\end{theorem}

\begin{proof}
The endpoint statements are definitions and the tail estimates.  Split the complete integral at $x$ and in the second part substitute $t=1-u$.  The elementary cases use the antiderivatives $t^a/a$ and $-(1-t)^b/b$.  All equalities are first established on rational approximations to $x$ and then transferred using the Lipschitz/error budget.
\end{proof}

\subsection{Recurrences and differentiation in the endpoint}

Integrating the derivative of $t^a(1-t)^b$ on $[0,x]$ gives
\[
 a\B_x(a,b+1)-b\B_x(a+1,b)=x^a(1-x)^b.
\]
Together with the incomplete Pascal identity
\[
 \B_x(a,b)=\B_x(a+1,b)+\B_x(a,b+1),
\]
this yields
\begin{align}
 \B_x(a+1,b)
 &=\frac{a}{a+b}\B_x(a,b)
   -\frac{x^a(1-x)^b}{a+b},
 \label{eq:inc-left-rec}\\
 \B_x(a,b+1)
 &=\frac{b}{a+b}\B_x(a,b)
   +\frac{x^a(1-x)^b}{a+b}.
 \label{eq:inc-right-rec}
\end{align}
After division by the corresponding complete-beta recurrence,
\begin{align}
 \I_x(a+1,b)
 &=\I_x(a,b)-\frac{x^a(1-x)^b}{a\B(a,b)},\\
 \I_x(a,b+1)
 &=\I_x(a,b)+\frac{x^a(1-x)^b}{b\B(a,b)}.
\end{align}

The finite fundamental theorem gives
\begin{align}
 \frac{d}{dx}\B_x(a,b)&=x^{a-1}(1-x)^{b-1},\\
 \frac{d}{dx}\I_x(a,b)&=
 \frac{x^{a-1}(1-x)^{b-1}}{\B(a,b)}.
\end{align}
These derivative statements require a local margin $0<d\le x\le1-d$ unless one uses one-sided endpoint certificates.

Integration by parts in $x$, again interpreted by finite telescoping, gives the source formulas
\begin{align*}
 \int\B_x(a,b)\,dx
 &=x\B_x(a,b)-\B_x(a+1,b)\\
 &=\left(x-\frac{a}{a+b}\right)\B_x(a,b)
   +\frac{x^a(1-x)^b}{a+b},\\
 \int\I_x(a,b)\,dx
 &=x\I_x(a,b)-\frac{a}{a+b}\I_x(a+1,b)\\
 &=\left(x-\frac{a}{a+b}\right)\I_x(a,b)
   +\frac{x^a(1-x)^b}{(a+b)\B(a,b)}.
\end{align*}
An indefinite integral here means the displayed antiderivative identity with a quadratic remainder, not a quotient by constants.

\subsection{Alternative integrals}

For $0<x<1$ with endpoint margins,
\begin{align}
 \B_x(a,b)
 &=2\int_0^{\arcsin\sqrt{x}}
   \sin^{2a-1}\theta\cos^{2b-1}\theta\,d\theta,
 \label{eq:inc-trig}\\
 &=\int_{(1-x)/x}^{\infty}
   \frac{s^{b-1}}{(1+s)^{a+b}}\,ds.
 \label{eq:inc-infinity}
\end{align}
The first uses $t=\sin^2\theta$ and a bisection name for the inverse on the monotone first quadrant.  The second uses $t=1/(1+s)$.  Its upper tail is bounded by $T^{-\Rea a}/\Rea a$.

The special identity
\begin{equation}
 \boxed{\B_x(a,-a)=\frac1a\left(\frac{x}{1-x}\right)^a}
 \label{eq:inc-special}
\end{equation}
follows by differentiating the right side and checking the value at $x=0$.  It is constructive for positive real $x<1$ and a specified positive-real power convention.  For complex parameters and transformed negative arguments, branch data must be included.

\subsection{Polynomial values at positive integers}

\begin{theorem}[Exact polynomial formula]\label{thm:inc-polynomial}
For positive integers $m,n$,
\begin{equation}
 \boxed{
 \B_x(m,n)=
 \sum_{j=0}^{n-1}(-1)^j\binom{n-1}{j}
 \frac{x^{m+j}}{m+j}.}
 \label{eq:inc-poly}
\end{equation}
It is a rational polynomial of degree $m+n-1$.
\end{theorem}

\begin{proof}
Use the finite binomial theorem
$(1-t)^{n-1}=\sum_{j=0}^{n-1}(-1)^j\binom{n-1}{j}t^j$
and integrate each monomial exactly.  This is a finite sum identity with no analytic remainder.
\end{proof}

The complement relation gives a second polynomial centered at $1$.  Multiplying by $(m+n-1)!/[(m-1)!(n-1)!]$ gives exact integer-coefficient formulas for the regularized function after common denominators are cleared.

\subsection{Hypergeometric representation}

For $0\le x\le1-\delta$,
\begin{align}
 \B_x(a,b)
 &=\sum_{n=0}^{\infty}
   \frac{(1-b)_n}{(a+n)n!}x^{a+n}
 \label{eq:inc-series}\\
 &=\frac{x^a}{a}\,
 {}_2F_1(a,1-b;a+1;x).
 \label{eq:inc-hyper}
\end{align}
The equality of the two displayed series uses
$(a)_n/(a+1)_n=a/(a+n)$, an exact finite-product cancellation.  The proof is the compact binomial-series proof of \cref{thm:beta-series}; the geometric margin $x\le1-\delta$ gives an especially simple cutoff.  At $x=1$ one instead uses the endpoint tail and the complete beta value.

A Lean development need not introduce a general hypergeometric-function object.  It may define the right side of \eqref{eq:inc-hyper} as this particular finite-sum name and prove the identity directly.

\subsection{The binomial cumulative-distribution identity}

For $0\le k<n$ and $0\le p\le1$, define the finite polynomial
\[
 F_{n,k}(p)=\sum_{j=0}^{k}\binom nj p^j(1-p)^{n-j}.
\]
No probability theory is needed.

\begin{theorem}[Binomial CDF identity]\label{thm:binomial-cdf}
\[
 \boxed{
 F_{n,k}(p)=\I_{1-p}(n-k,k+1)
 =1-\I_p(k+1,n-k).}
\]
\end{theorem}

\begin{proof}
The two regularized beta expressions are equal by complementarity.  Differentiate the finite polynomial.  Using
$j\binom nj=n\binom{n-1}{j-1}$ and
$(n-j)\binom nj=n\binom{n-1}{j}$, adjacent terms telescope:
\[
 F'_{n,k}(p)
 =-\frac{n!}{k!(n-k-1)!}p^k(1-p)^{n-k-1}.
\]
By the factorial formula for $\B(k+1,n-k)$, this is exactly the derivative of
$1-\I_p(k+1,n-k)$.  Both sides equal $1$ at $p=0$.  The finite fundamental theorem proves equality.  Alternatively, induction on $n$ using Pascal's identity stays entirely in polynomial arithmetic.  The omitted boundary case $k=n$ has $F_{n,n}(p)=1$ by the binomial theorem; it is handled separately because the regularized parameter $n-k$ would be zero.
\end{proof}

\subsection{Continued fraction: exact finite target and remaining lemma}

The source gives
\[
 \B_x(a,b)=
 \frac{x^a(1-x)^b}
 {a\left(1+\cfrac{d_1}{1+\cfrac{d_2}{1+\cdots}}\right)}
\]
with
\[
 d_{2m+1}=
 -\frac{(a+m)(a+b+m)x}{(a+2m)(a+2m+1)},
 \qquad
 d_{2m}=
 \frac{m(b-m)x}{(a+2m-1)(a+2m)}.
\]
Define continuants by
\[
 P_{-1}=1,\ P_0=1,\quad
 Q_{-1}=0,\ Q_0=1,\quad
 P_n=P_{n-1}+d_nP_{n-2},\quad
 Q_n=Q_{n-1}+d_nQ_{n-2}.
\]
The $n$th finite continued fraction is $P_n/Q_n$, and its determinant identity
\[
 P_nQ_{n-1}-P_{n-1}Q_n
 =(-1)^{n-1}d_1d_2\cdots d_n
\]
is a one-line induction, fully internal to iPRA.

The remaining analytic theorem is a positive-remainder representation comparing
\[
 \frac{x^a(1-x)^bQ_n}{aP_n}
\]
with $\B_x(a,b)$.  Classical proofs obtain it from Gauss contiguous relations for ${}_2F_1$; the signs of the remainder yield the stated $4m/4m+1$ lower and $4m+2/4m+3$ upper convergents.  A complete iPRA proof should:
\begin{enumerate}
\item express every contiguous relation as a finite coefficient identity for truncated hypergeometric sums;
\item attach the geometric tail from $x\le1-\delta$;
\item derive the remainder as an integral with a nonnegative rational-mesh integrand for $a,b>0$;
\item use the complement relation when $x$ is near $1$.
\end{enumerate}
The continuant algebra and formal target are complete; the general alternating-remainder lemma is classified \statusC{} in this packet.
